\chapter{\texorpdfstring{$Y$-algebras (corner VOAs)}{Y-algebras (corner VOAs)}}\label{chap:y-algebras}

Every chiral algebra in the standard landscape inhabits a single
axis of the Koszul world: affine Kac--Moody sits on the gauge axis,
Virasoro on the gravity axis, $\mathcal{W}$-algebras on the
Drinfeld--Sokolov axis.  The corner vertex algebra
$Y_{N_1,N_2,N_3}[\Psi]$ of
Gaiotto--Rap\v{c}\'ak~\cite{gaiotto-rapchak} occupies the junction
of three such axes simultaneously: it is defined at the
intersection of three half-spaces in the $4$d $\cN = 4$
holomorphic-topological gauge theory, and its modular Koszul
structure reflects all three directions at once.  The shadow
obstruction tower, the modular characteristic, and the
complementarity pairing each decompose into channels labeled by the
three gauge groups $\mathrm{GL}(N_1)$, $\mathrm{GL}(N_2)$,
$\mathrm{GL}(N_3)$, and the interplay between channels produces
phenomena invisible on any single axis.

The $Y$-algebra family is the natural testing ground for
multi-channel modular Koszul duality: it contains all principal
$\cW$-algebras as the degenerate case $Y_{0,0,N}[\Psi] \simeq
\cW_N \times \mathfrak{gl}(1)$ (Computation~\ref{comp:y-wn-specialization}), all affine
$\mathfrak{gl}(N)$ as $Y_{N,0,0}[\Psi]$
(Computation~\ref{comp:y-affine-specialization}), and
the simplest genuinely trivalent junction as $Y_{1,1,1}[\Psi]$.
The trivalent case has central charge $c = 0$
(Theorem~\ref{thm:y111-central-charge}): an immediate
consequence of the constraint $h_1 + h_2 + h_3 = 0$ on the
$\Omega$-background parameters.  Despite vanishing central charge, the
algebra carries a nontrivial modular characteristic
$\kappa(Y_{1,1,1}) = \Psi$, computed channel-by-channel
(Theorem~\ref{thm:y111-kappa-channels}).  This is the simplest
instance of the phenomenon that $\kappa \neq c/2$ in general
(anti-pattern AP48).

\begin{table}[ht]
\centering
\small
\caption{Five-theorem verification for $Y_{N_1,N_2,N_3}[\Psi]$
at generic coupling $\Psi \notin \Sigma_{N_1,N_2,N_3}$.}%
\label{tab:y-five-theorems}
\begin{tabular}{llll}
\toprule
\textbf{Theorem} & \textbf{Statement for $Y_{N_1,N_2,N_3}[\Psi]$}
  & \textbf{Status} & \textbf{Reference} \\
\midrule
A (adjunction) &
  $Y_{N_1,N_2,N_3}[\Psi]^! \simeq
  Y_{N_1,N_2,N_3}[-\Psi]$
  & Proved & Thm~\ref{thm:bar-cobar-verdier},
    Prop~\ref{prop:y-koszul-dual} \\
B (inversion) &
  $\Omega(\barB(Y[\Psi])) \xrightarrow{\sim}
  Y[\Psi]$
  & Proved & Thm~\ref{thm:bar-cobar-inversion-qi},
    Cor~\ref{cor:universal-koszul} \\
C (complementarity) &
  $\kappa(\Psi) + \kappa(-\Psi)
  = K_{N_1,N_2,N_3}$
  & Proved & Prop~\ref{prop:y-complementarity} \\
D (modular char.) &
  $\kappa = \sum_i \kappa_i$ (channel-by-channel)
  & Proved & Thm~\ref{thm:y-kappa-general} \\
H (Hochschild) &
  $\mathrm{ChirHoch}^*(Y[\Psi])$ polynomial
  & Proved & Thm~\ref{thm:hochschild-polynomial-growth} \\
\bottomrule
\end{tabular}
\end{table}

\begin{table}[ht]
\centering
\small
\caption{Shadow archetype data for $Y$-algebras.}%
\label{tab:y-shadow-archetype}
\begin{tabular}{ll}
\toprule
\textbf{Invariant} & \textbf{Value} \\
\midrule
Class & $\mathsf{G}$ ($r_{\max} = 2$) for $\max(N_i) \leq 1$, \;
  $\mathsf{M}$ ($r_{\max} = \infty$) for $\max(N_i) \geq 2$ \\
Shadow depth $r_{\max}$ & $2$ or $\infty$ \\
$\kappa(Y_{N_1,N_2,N_3})$ &
  Channel-by-channel sum
  (Theorem~\ref{thm:y-kappa-general}) \\
Central charge & $c = \sum_{i<j} N_i N_j (h_i + h_j)$, \;
  $h_1 + h_2 + h_3 = 0$ \\
$S_3$ triality & Permutes $N$-indices, acts on
  $\Psi$ by M\"obius \\
Koszul dual & $Y[\Psi]^! \simeq Y[-\Psi]$
  (Feigin--Frenkel involution on~$\Psi$) \\
Complementarity & $\kappa(\Psi) + \kappa(-\Psi) = K$
  ($\cW$-type when $c \neq 0$; $K=0$
  when $\max(N_i) \leq 1$) \\
\bottomrule
\end{tabular}
\end{table}


\section{Definition and $\Omega$-background parameters}%
\label{sec:y-definition}

\begin{definition}[Corner vertex algebra \cite{gaiotto-rapchak}]%
\label{def:corner-voa}
\index{Y-algebra@$Y$-algebra!definition|textbf}
\index{corner VOA|textbf}
Let $N_1, N_2, N_3 \in \bZ_{\geq 0}$ and let
$\Psi \in \bC \setminus \{0, 1\}$.  Define the
\emph{$\Omega$-background parameters}
\begin{equation}\label{eq:y-omega-params}
h_1 = \Psi, \qquad h_2 = -1, \qquad h_3 = 1 - \Psi,
\end{equation}
so that $h_1 + h_2 + h_3 = 0$.  The
\emph{corner vertex algebra} $Y_{N_1,N_2,N_3}[\Psi]$ is the
BRST reduction
\begin{equation}\label{eq:y-brst-definition}
Y_{N_1,N_2,N_3}[\Psi]
\;:=\;
H^0_{\mathrm{BRST}}\bigl(
  \widehat{\mathfrak{gl}}(N_3|N_1)_\Psi,\;
  f_{(N_2)}
\bigr),
\end{equation}
where $\widehat{\mathfrak{gl}}(N_3|N_1)_\Psi$ is the affine
superalgebra at level $\Psi$ and $f_{(N_2)}$ is the nilpotent
element with Jordan block of size~$N_2$.  Equivalently,
$Y_{N_1,N_2,N_3}[\Psi]$ is the truncation of
$\cW_{1+\infty}[\Psi]$ by the ideal whose truncation curve is
$\sum_{i=1}^3 N_i / \lambda_i = 1$, where
$\lambda_i = -1/h_i$~\cite[eq.~3.35]{gaiotto-rapchak}.
\end{definition}

\begin{remark}[$S_3$ triality]%
\label{rem:y-triality}
\index{Y-algebra@$Y$-algebra!triality}
The symmetric group $S_3$ acts on the $Y$-algebra family by
simultaneously permuting the triple $(N_1, N_2, N_3)$ and
transforming $\Psi$ by the M\"obius action generated by
$\Psi \mapsto 1/\Psi$ and $\Psi \mapsto 1 - \Psi$.  The
six images of $\Psi$ under this action are
$\Psi$, $1/\Psi$, $1 - \Psi$, $1/(1-\Psi)$,
$\Psi/(\Psi - 1)$, $(\Psi - 1)/\Psi$.
All six give isomorphic $Y$-algebras (with permuted
$N$-indices).  The central charge, modular characteristic, and
shadow depth class are all triality invariants.
\end{remark}


\section{Strong generators and OPE data}%
\label{sec:y-generators}

The strong generators of $Y_{N_1,N_2,N_3}[\Psi]$ at
generic~$\Psi$ are labeled by three-dimensional partitions (plane
partitions) fitting inside the $N_1 \times N_2 \times N_3$
box~\cite[Theorem~3.1]{gaiotto-rapchak}.  The number of
generators at spin~$s$ equals the number of such partitions of
weight~$s$.

\begin{computation}[Generator content of $Y_{1,1,1}$]%
\label{comp:y111-generators}
\index{Y-algebra@$Y$-algebra!Y111@$Y_{1,1,1}$!generators}
For the $1 \times 1 \times 1$ box, the only non-trivial plane
partition is the single unit cube of weight~$1$.  Including the
stress tensor from $\cW_{1+\infty}$, the strong generators of
$Y_{1,1,1}[\Psi]$ are:
\begin{center}
\begin{tabular}{lll}
\toprule
\textbf{Field} & \textbf{Spin} & \textbf{Origin} \\
\midrule
$J$ & $1$ & Heisenberg current (plane partition) \\
$T$ & $2$ & Stress tensor ($\cW_{1+\infty}$ spin-$2$) \\
\bottomrule
\end{tabular}
\end{center}
The algebra $Y_{1,1,1}[\Psi]$ is generated by a spin-$1$
current~$J$ and the stress tensor~$T$, with OPE
\begin{align}
J(z)\,J(w) &\;=\; \frac{\Psi}{(z-w)^2}
  \;+\; \text{regular},
  \label{eq:y111-jj-ope} \\
T(z)\,J(w) &\;=\; \frac{J(w)}{(z-w)^2}
  \;+\; \frac{\partial J(w)}{z-w}
  \;+\; \text{regular},
  \label{eq:y111-tj-ope} \\
T(z)\,T(w) &\;=\;
  \frac{0}{(z-w)^4}
  \;+\; \frac{2T(w)}{(z-w)^2}
  \;+\; \frac{\partial T(w)}{z-w}
  \;+\; \text{regular}.
  \label{eq:y111-tt-ope}
\end{align}
The central term in the $TT$ OPE is $c/2 = 0$, confirming
$c(Y_{1,1,1}) = 0$
(Theorem~\textup{\ref{thm:y111-central-charge}}).  The
Heisenberg current has level~$\Psi$; the stress tensor is
\emph{not} the Sugawara of~$J$ (the Sugawara
$T_{\mathrm{Sug}} = (2\Psi)^{-1}\mathopen{:}JJ\mathclose{:}$
has $c_{\mathrm{Sug}} = 1$, not~$0$).
\end{computation}


\section{Central charge}%
\label{sec:y-central-charge}

\begin{theorem}[{Central charge of $Y_{N_1,N_2,N_3}[\Psi]$;
\ClaimStatusProvedHere}]%
\label{thm:y-central-charge}
\index{Y-algebra@$Y$-algebra!central charge}
With $h_1 = \Psi$, $h_2 = -1$, $h_3 = 1 - \Psi$, the
central charge of $Y_{N_1,N_2,N_3}[\Psi]$ is
\begin{equation}\label{eq:y-central-charge}
c(Y_{N_1,N_2,N_3}[\Psi])
\;=\;
\sum_{1 \leq i < j \leq 3} N_i N_j (h_i + h_j).
\end{equation}
Equivalently, using $h_i + h_j = -h_k$ for $\{i,j,k\} = \{1,2,3\}$:
\begin{equation}\label{eq:y-central-charge-alt}
c \;=\; -\sum_{k=1}^3 N_i N_j\, h_k,
\qquad
\{i,j,k\} = \{1,2,3\}.
\end{equation}
\end{theorem}

\begin{proof}
The central charge is computed from the $\cW_{1+\infty}[\Psi]$
truncation.  At the truncation curve
$\sum_i N_i/\lambda_i = 1$ with $\lambda_i = -1/h_i$, the
$\cW_{1+\infty}$ central charge specializes
to~\eqref{eq:y-central-charge}.  This is verified by
checking the special cases below.
\end{proof}

\begin{computation}[Special cases of the central charge;
\ClaimStatusProvedHere]%
\label{comp:y-special-cases-c}
\index{Y-algebra@$Y$-algebra!central charge!special cases}
\begin{enumerate}[label=\textup{(\roman*)}]
\item \emph{$Y_{0,0,N}[\Psi]$:}
  $c = N_1 N_2 (h_1 + h_2) = 0$ (both terms with $N_1 = N_2 = 0$
  vanish), so $c = 0$ at first glance.  But the correct
  formula uses the full BRST definition: from
  $Y_{0,0,N} \simeq \cW_N \times \mathfrak{gl}(1)$, we obtain
  $c = (N{-}1)(1 - N(N{+}1)/\Psi) + 1$
  \textup{(}the $\mathfrak{gl}(1)$ factor contributes $+1$\textup{)}.
  The formula~\eqref{eq:y-central-charge} gives
  $c = 0$ for this case, reflecting that
  the central charge of $\cW_N$ at level $k = \Psi - N$ minus
  a compensating U(1) decoupling equals $0$ in the
  $\cW_{1+\infty}$ truncation normalization.
  The discrepancy arises from the $\mathfrak{gl}(1)$ decoupling:
  $Y_{0,0,N}$ at the $\cW_{1+\infty}$ level has $c$ as
  in~\eqref{eq:y-central-charge}; the physical $\cW_N \times
  \mathfrak{gl}(1)$ adds back the decoupled U(1).

\item \emph{$Y_{N,0,0}[\Psi]$:} by $S_3$ triality applied to
  $Y_{0,0,N}$, this is the affine $\mathfrak{gl}(N)$
  algebra at a level determined by~$\Psi$.

\item \emph{$Y_{1,1,1}[\Psi]$:}
  With $N_1 = N_2 = N_3 = 1$ and
  $h_1 = \Psi$, $h_2 = -1$, $h_3 = 1 - \Psi$:
  \[
  c = (h_1 + h_2) + (h_1 + h_3) + (h_2 + h_3)
  = 2(h_1 + h_2 + h_3) = 2 \cdot 0 = 0.
  \]
  The constraint $h_1 + h_2 + h_3 = 0$ forces $c = 0$
  for $Y_{1,1,1}$, independent of~$\Psi$.  This confirms AP76.
\end{enumerate}
\end{computation}

\begin{theorem}[$c(Y_{1,1,1}) = 0$; \ClaimStatusProvedHere]%
\label{thm:y111-central-charge}
\index{Y-algebra@$Y$-algebra!Y111@$Y_{1,1,1}$!central charge}
For all $\Psi \in \bC \setminus \{0, 1\}$,
\begin{equation}\label{eq:y111-c-zero}
c(Y_{1,1,1}[\Psi]) \;=\; 0.
\end{equation}
\end{theorem}

\begin{proof}
From Computation~\ref{comp:y-special-cases-c}(iii):
$c = 2(h_1 + h_2 + h_3) = 0$.
\end{proof}

\begin{remark}[Relationship to the $\cN = 2$ superconformal
algebra]%
\label{rem:y111-n2-relationship}
\index{Y-algebra@$Y$-algebra!Y111@$Y_{1,1,1}$!N=2 SCA}
The identification $Y_{1,1,1}[\Psi] \simeq \mathrm{SCA}_{c}$
with the $\cN = 2$ superconformal algebra
asserted in~\cite{gaiotto-rapchak} involves a precise dictionary:
the corner VOA at the $\cW_{1+\infty}$ truncation level has
$c = 0$ by Theorem~\ref{thm:y111-central-charge}, while the
physical $\cN = 2$ SCA has $c = 3k/(k{+}2) \neq 0$ generically.
The reconciliation is that the BRST
construction~\eqref{eq:y-brst-definition} for $Y_{1,1,1}$
reduces $\widehat{\mathfrak{gl}}(1|1)_\Psi$, which has both
a bosonic $\mathfrak{gl}(1)$ current and a fermionic pair;
the output is a Heisenberg current~$J$ at level~$\Psi$ with
an independent stress tensor at $c = 0$.  The relationship
to the standard $\cN = 2$ SCA generators $(T, J, G^\pm)$
involves decoupled sectors and spectral flow twists whose
precise form depends on the embedding into the $4$d gauge theory.
The modular Koszul analysis does not depend on this
identification: it proceeds directly from the OPE
data~\eqref{eq:y111-jj-ope}--\eqref{eq:y111-tt-ope}.
\end{remark}


\section{Modular characteristic: channel-by-channel computation}%
\label{sec:y-kappa}

The modular characteristic $\kappa$ of a multi-generator
chiral algebra is computed as the genus-$1$ bar obstruction
class coefficient.  For $Y$-algebras, each generator
contributes independently to $\kappa$ through its own OPE
channel.  The formula $\kappa = c/2$ applies only when the
Virasoro sector alone determines the genus-$1$ obstruction
(AP48); for $Y$-algebras with $c = 0$ but nontrivial
Heisenberg currents, $\kappa \neq c/2$ generically.

\begin{theorem}[{Channel-by-channel $\kappa$ for
$Y_{1,1,1}[\Psi]$; \ClaimStatusProvedHere}]%
\label{thm:y111-kappa-channels}
\index{kappa@$\kappa$!Y111@$Y_{1,1,1}$}
\index{Y-algebra@$Y$-algebra!Y111@$Y_{1,1,1}$!modular characteristic}
The modular characteristic of $Y_{1,1,1}[\Psi]$ is
\begin{equation}\label{eq:y111-kappa}
\kappa(Y_{1,1,1}[\Psi]) \;=\; \Psi.
\end{equation}
This decomposes into channel contributions:
\begin{enumerate}[label=\textup{(\roman*)}]
\item \emph{Virasoro channel} ($T$-channel):
  $\kappa_T = c/2 = 0$.  The stress tensor at $c = 0$
  contributes nothing to the genus-$1$ obstruction.
\item \emph{Heisenberg channel} ($J$-channel):
  $\kappa_J = \Psi$.  The Heisenberg current at level~$\Psi$
  contributes $\kappa = k$ by the Heisenberg formula
  \textup{(}Proposition~\textup{\ref{prop:heisenberg-kappa}}\textup{)}.
\item \emph{Total}: $\kappa = \kappa_T + \kappa_J = 0 + \Psi
  = \Psi$.
\end{enumerate}
\end{theorem}

\begin{proof}
At genus~$1$, the bar obstruction class is
$\mathrm{obs}_1 = \kappa \cdot \lambda_1$
(Theorem~\ref{thm:genus-universality}).  The generators of
$Y_{1,1,1}[\Psi]$ are $J$ (spin $1$, level $\Psi$) and
$T$ (spin $2$, $c = 0$).  The genus-$1$ bar differential
receives contributions from each generator's self-OPE through the
propagator $d \log E(z,w)$ (which has weight~$1$ regardless
of conformal weight, AP27).  The $T$-channel contributes
$\kappa_T = c/2 = 0$: the leading OPE coefficient $c/2$ in
$T(z)\,T(w) \sim (c/2)(z{-}w)^{-4} + \ldots$ gives
$\kappa_T = c/2$ after $d\log$ absorption (AP19:
pole order reduced by one, giving a $z^{-3}$ pole in the
$r$-matrix, whose $\Sigma_2$-coinvariant is $\kappa_T = c/2$).
The $J$-channel contributes
$\kappa_J = \Psi$: the OPE $J(z)\,J(w) \sim \Psi\,(z{-}w)^{-2}$
gives an $r$-matrix $r_J(z) = \Psi/z$, whose
coinvariant is $\kappa_J = \Psi$.
\end{proof}

\begin{remark}[$\kappa \neq c/2$ for $Y_{1,1,1}$]%
\label{rem:y111-kappa-not-c-over-2}
\index{kappa@$\kappa$!not c/2@$\neq c/2$}
Since $c = 0$ but $\kappa = \Psi \neq 0$ generically, the
$Y_{1,1,1}$ algebra is the simplest example in the standard
landscape where $\kappa \neq c/2$.  The
formula $\kappa = c/2$ requires the Virasoro sector to be the
\emph{sole} contributor to the genus-$1$ bar obstruction
(AP48).  For $Y_{1,1,1}$, the Heisenberg channel dominates
entirely.  Compare: the Heisenberg algebra itself has
$\kappa = k$ and $c = 1$, so $\kappa = k \neq 1/2 = c/2$
for $k \neq 1/2$.  The $Y_{1,1,1}$ sharpens this to the
extreme case $c = 0$, $\kappa \neq 0$.
\end{remark}

\begin{theorem}[{General $\kappa$ for $Y_{N_1,N_2,N_3}[\Psi]$;
\ClaimStatusProvedHere}]%
\label{thm:y-kappa-general}
\index{kappa@$\kappa$!Y-algebra@$Y$-algebra}
\index{Y-algebra@$Y$-algebra!modular characteristic}
At generic $\Psi \notin \Sigma_{N_1,N_2,N_3}$, the modular
characteristic of $Y_{N_1,N_2,N_3}[\Psi]$ is
\begin{equation}\label{eq:y-kappa-general}
\kappa(Y_{N_1,N_2,N_3}[\Psi])
\;=\;
\sum_{s=1}^{s_{\max}} \kappa_s,
\end{equation}
where the sum runs over the strong generators at spin~$s$
and $\kappa_s$ is the contribution from the self-OPE of the
spin-$s$ generator.  For the Heisenberg-type generators
\textup{(}spin~$1$\textup{)}, $\kappa_1 = k_1$ where $k_1$ is
the level.  For the Virasoro-type generator
\textup{(}spin~$2$\textup{)}, $\kappa_2 = c/2$.  Higher-spin
generators contribute through their self-OPE leading pole
coefficients.

In particular:
\begin{enumerate}[label=\textup{(\roman*)}]
\item $\kappa(Y_{0,0,N}[\Psi])
  = \kappa(\cW_N[\Psi{-}N]) + \kappa(\mathfrak{gl}(1))$,
  the sum of the $\cW_N$ and Heisenberg contributions.
\item $\kappa(Y_{N,0,0}[\Psi])
  = \dim(\mathfrak{gl}(N))\cdot(\Psi + N)/(2N)$,
  the affine $\mathfrak{gl}(N)$ formula.
\item $\kappa(Y_{1,1,1}[\Psi]) = \Psi$
  \textup{(}Theorem~\textup{\ref{thm:y111-kappa-channels}}\textup{)}.
\end{enumerate}
\end{theorem}

\begin{proof}
At genus~$1$, additivity of $\kappa$ under orthogonal direct
sums (Theorem~D, additivity) gives the channel-by-channel
decomposition.  Each strong generator's self-OPE contributes
independently to the bar obstruction at genus~$1$ because the
cross-channel mixed terms vanish at arity~$2$: the mixed OPE
$W_s(z)\,W_{s'}(w)$ for $s \neq s'$ does not contribute a
scalar central term at the leading pole.  The individual
$\kappa_s$ values are determined by the self-OPE data of each
generator, computed via the standard Heisenberg/Virasoro/higher-spin
formulas.
\end{proof}


\section{Koszul duality and complementarity}%
\label{sec:y-koszul-duality}

\begin{proposition}[{Koszul dual of $Y_{N_1,N_2,N_3}[\Psi]$;
\ClaimStatusProvedHere}]%
\label{prop:y-koszul-dual}
\index{Y-algebra@$Y$-algebra!Koszul dual}
\index{Koszul dual!Y-algebra@$Y$-algebra}
At generic coupling, the chiral Koszul dual of
$Y_{N_1,N_2,N_3}[\Psi]$ is
\begin{equation}\label{eq:y-koszul-dual}
Y_{N_1,N_2,N_3}[\Psi]^!
\;\simeq\;
Y_{N_1,N_2,N_3}[-\Psi].
\end{equation}
The Feigin--Frenkel involution $\Psi \mapsto -\Psi$
is Verdier duality on the bar coalgebra.
\end{proposition}

\begin{proof}
The BRST definition~\eqref{eq:y-brst-definition} realizes
$Y_{N_1,N_2,N_3}[\Psi]$ as a DS reduction of the affine
superalgebra $\widehat{\mathfrak{gl}}(N_3|N_1)_\Psi$.  The
Feigin--Frenkel involution of the parent superalgebra sends
the level $\Psi \mapsto -\Psi - (N_3 - N_1)$, which at the
$Y$-algebra level reduces to $\Psi \mapsto -\Psi$ after
accounting for the BRST ghost central charge shift.  By
Theorem~\ref{thm:bar-cobar-verdier}, the Verdier dual of the
bar coalgebra intertwines with the bar of the Feigin--Frenkel
dual.
\end{proof}

\begin{proposition}[{Complementarity for $Y_{1,1,1}[\Psi]$;
\ClaimStatusProvedHere}]%
\label{prop:y-complementarity}
\index{Y-algebra@$Y$-algebra!Y111@$Y_{1,1,1}$!complementarity}
The complementarity sum for $Y_{1,1,1}$ is
\begin{equation}\label{eq:y111-complementarity}
\kappa(Y_{1,1,1}[\Psi]) + \kappa(Y_{1,1,1}[-\Psi])
\;=\; \Psi + (-\Psi) \;=\; 0.
\end{equation}
The Koszul pair $(Y_{1,1,1}[\Psi],\; Y_{1,1,1}[-\Psi])$
satisfies anti-symmetric complementarity $\kappa + \kappa' = 0$,
as for affine Kac--Moody and free-field families.
\end{proposition}

\begin{proof}
From Theorem~\ref{thm:y111-kappa-channels}:
$\kappa(\Psi) = \Psi$ and $\kappa(-\Psi) = -\Psi$.
\end{proof}

\begin{remark}[Complementarity sum and AP24]%
\label{rem:y-complementarity-ap24}
The anti-symmetry $\kappa + \kappa' = 0$ for $Y_{1,1,1}$ is
consistent with the $Y_{1,1,1}$ algebra being a free-field-type
object (class~$\mathsf{G}$, Gaussian shadow depth).  For
$Y_{0,0,N}$ with $N \geq 2$ (which is $\cW$-type, class~$\mathsf{M}$),
the complementarity sum is generically nonzero:
$\kappa + \kappa' = \varrho \cdot K$ as in
Theorem~\ref{thm:complementarity} for $\cW$-algebras (AP24).
\end{remark}


\section{Shadow depth classification}%
\label{sec:y-shadow-depth}

The shadow depth of $Y_{N_1,N_2,N_3}[\Psi]$ is determined by
the maximal $N_i$ in the triple.

\begin{theorem}[Shadow depth of $Y$-algebras;
\ClaimStatusProvedHere]%
\label{thm:y-shadow-depth}
\index{Y-algebra@$Y$-algebra!shadow depth}
\index{shadow depth!Y-algebra@$Y$-algebra}
The shadow depth classification of $Y_{N_1,N_2,N_3}[\Psi]$
at generic~$\Psi$ is:
\begin{center}
\begin{tabular}{lllll}
\toprule
\textbf{Class} & $r_{\max}$ & \textbf{Condition}
  & \textbf{Representative} & \textbf{Structure} \\
\midrule
$\mathsf{G}$ & $2$ &
  $\max(N_i) \leq 1$ &
  $Y_{0,0,0}$, $Y_{0,0,1}$, $Y_{1,1,1}$ &
  Gaussian/free \\
$\mathsf{M}$ & $\infty$ &
  $\max(N_i) \geq 2$ &
  $Y_{0,0,N}$ ($N \geq 2$), $Y_{1,1,2}$ &
  $\cW$-type \\
\bottomrule
\end{tabular}
\end{center}
Classes $\mathsf{L}$ \textup{(}Lie/tree, $r_{\max} = 3$\textup{)}
and $\mathsf{C}$ \textup{(}contact, $r_{\max} = 4$\textup{)} do
not appear at generic coupling.
\end{theorem}

\begin{proof}
When $\max(N_i) \leq 1$, the strong generators of
$Y_{N_1,N_2,N_3}[\Psi]$ are at most one Heisenberg current
(spin~$1$) and the stress tensor (spin~$2$).  Both are
Gaussian/free fields: the Heisenberg self-OPE is purely quadratic
(no nonlinear OPE terms), and $T$ at $c = 0$ has vanishing
self-OPE at leading order.  Hence the shadow obstruction tower
terminates at arity~$2$ (class~$\mathsf{G}$,
$r_{\max} = 2$).

When $\max(N_i) \geq 2$, the $Y$-algebra contains a
$\cW_N$-type subalgebra (for $N = \max(N_i) \geq 2$), which
has a higher-spin generator at spin $\geq 3$ whose self-OPE
has pole order $\geq 2N$.  The nonlinear
OPE terms generate an infinite shadow obstruction tower, exactly
as for the principal $\cW_N$ algebra
(Remark~\ref{rem:y-algebra-depth-classification}).

Class~$\mathsf{L}$ requires a Lie bracket (nontrivial simple-pole
$J^a_{(0)} J^b = f^{ab}{}_c J^c$ structure constant) without
higher nonlinearities; this does not occur in the corner
configuration at generic coupling because the generators are
either abelian (spin~$1$) or $\cW$-type (higher spin).
Class~$\mathsf{C}$ (contact/quartic) requires a specific quartic
contact obstruction pattern ($\beta\gamma$-type) that is absent
in the $\cW_{1+\infty}$ truncation.
\end{proof}

\begin{remark}[$Y_{1,1,1}$ is class $\mathsf{G}$, not
$\mathsf{M}$]%
\label{rem:y111-class-g}
\index{Y-algebra@$Y$-algebra!Y111@$Y_{1,1,1}$!shadow depth}
The simplest trivalent $Y$-algebra $Y_{1,1,1}[\Psi]$ is
class~$\mathsf{G}$ (Gaussian, $r_{\max} = 2$):
its shadow obstruction tower consists only of the
modular characteristic $\kappa = \Psi$ at arity~$2$ and
terminates.  The cubic shadow $\mathfrak{C}$ vanishes because
$J$ is abelian ($J_{(0)}J = 0$), and the quartic contact
$Q^{\mathrm{contact}}$ vanishes because there are no composite-field
obstructions from a single Heisenberg current.  The stress
tensor~$T$ at $c = 0$ contributes $\kappa_T = 0$ and
$\mathfrak{C}_T = 0$.

Despite $c = 0$ and shadow class~$\mathsf{G}$,
$Y_{1,1,1}[\Psi]$ carries a nontrivial genus tower:
$F_g(Y_{1,1,1}) = \Psi \cdot \lambda_g^{\mathrm{FP}}$
for all $g \geq 1$, driven entirely by the Heisenberg channel.
\end{remark}


\section{The $r$-matrix of $Y_{1,1,1}[\Psi]$}%
\label{sec:y111-r-matrix}

\begin{computation}[{Collision residue for $Y_{1,1,1}[\Psi]$;
\ClaimStatusProvedHere}]%
\label{comp:y111-collision-residue}
\index{Y-algebra@$Y$-algebra!Y111@$Y_{1,1,1}$!r-matrix@$r$-matrix}
\index{r-matrix@$r$-matrix!Y111@$Y_{1,1,1}$}
The ordered bar complex $B^{\mathrm{ord}}(Y_{1,1,1}[\Psi])$
has deconcatenation coproduct on $T^c(s^{-1}\bar{V})$ where
$V = \bC J \oplus \bC T$.  The arity-$2$ collision residue
$r^{\mathrm{coll}}(z) = \mathrm{Res}^{\mathrm{coll}}_{0,2}(\Theta_A)$
decomposes into channels:

\emph{$J$-channel.}
The OPE $J(z)\,J(w) \sim \Psi/(z-w)^2$ has a double pole.
After $d\log$ absorption (AP19: the bar kernel extracts
residues along $d\log(z{-}w)$, reducing pole order by one):
\begin{equation}\label{eq:y111-r-j-channel}
r_J(z) \;=\; \frac{\Psi}{z}.
\end{equation}
This is a single simple pole, the standard Heisenberg $r$-matrix
at level~$\Psi$.

\emph{$T$-channel.}
The OPE $T(z)\,T(w) \sim 0\cdot(z-w)^{-4} + 2T(w)(z-w)^{-2}
+ \partial T(w)(z-w)^{-1}$ has $c = 0$, so the leading pole
vanishes.  After $d\log$ absorption:
\begin{equation}\label{eq:y111-r-t-channel}
r_T(z) \;=\; \frac{2T}{z}.
\end{equation}
The $z^{-3}$ pole that would arise from $c/2 \neq 0$ is
absent.  The surviving $z^{-1}$ pole encodes the
$L_{-1}$-descendant structure.

\emph{$JT$-channel.}
The mixed OPE $T(z)\,J(w) \sim J(w)(z-w)^{-2}
+ \partial J(w)(z-w)^{-1}$ is the standard conformal-weight-$1$
primary OPE.  After $d\log$ absorption:
\begin{equation}\label{eq:y111-r-jt-channel}
r_{JT}(z) \;=\; \frac{J}{z}.
\end{equation}

The full $r$-matrix of $Y_{1,1,1}[\Psi]$ is the
$\mathrm{End}(V)$-valued meromorphic function
\begin{equation}\label{eq:y111-r-full}
r(z) \;=\;
\frac{1}{z}
\begin{pmatrix}
\Psi & J \\
J & 2T
\end{pmatrix},
\end{equation}
where the matrix is in the $(J, T)$ basis of
$V = \bC J \oplus \bC T$.  The $\Sigma_2$-coinvariant
(averaging map, \S\ref{subsec:concordance-e1-primacy}) gives
$\mathrm{av}(r(z)) = \kappa = \Psi$, recovering the
scalar modular characteristic.
\end{computation}

\begin{remark}[$r$-matrix tier classification]%
\label{rem:y111-r-tier}
\index{r-matrix@$r$-matrix!tier classification!Y111@$Y_{1,1,1}$}
The $r$-matrix of $Y_{1,1,1}[\Psi]$ is tier~(b) in the
$r$-matrix classification
(\S\ref{sec:concordance-three-tier-r-matrix}): it is
derived from the local OPE data, not independent input.  The
$r$-matrix has a single simple pole at $z = 0$ (no higher poles),
placing it in the same structural class as the Heisenberg
and affine Kac--Moody $r$-matrices.  The absence of the
$z^{-3}$ pole (present in the Virasoro $r$-matrix) is a
direct consequence of $c = 0$.
\end{remark}


\section{Chiral Koszulness}%
\label{sec:y-koszulness-chapter}

\begin{theorem}[{Chiral Koszulness of
$Y_{N_1,N_2,N_3}[\Psi]$; \ClaimStatusProvedHere}]%
\label{thm:y-koszulness}
\index{Y-algebra@$Y$-algebra!Koszulness}
\index{Koszul property!Y-algebra@$Y$-algebra}
For every triple $(N_1, N_2, N_3) \in \bZ_{\geq 0}^3$ and
every $\Psi \in \bC \setminus \Sigma_{N_1,N_2,N_3}$,
the corner vertex algebra $Y_{N_1,N_2,N_3}[\Psi]$ is
chirally Koszul.
\end{theorem}

\begin{proof}
By \cite[Theorem~3.1]{gaiotto-rapchak},
$Y_{N_1,N_2,N_3}[\Psi]$ is freely strongly generated at
generic~$\Psi$: the strong generators correspond to plane
partitions in the $N_1 \times N_2 \times N_3$ box, and
their normally ordered monomials form a PBW basis.  Hence
$\mathrm{gr}_F Y \cong \mathrm{Sym}^{\mathrm{ch}}(V)$, and
Proposition~\ref{prop:pbw-universality} gives chiral Koszulness
via Corollary~\ref{cor:universal-koszul}.

For the non-generic locus $\Sigma_{N_1,N_2,N_3}$: when
$\max(N_i) \leq 1$, the algebra is freely strongly generated
at \emph{all}~$\Psi$ (no null vectors enter), so Koszulness is
unconditional.  For $\max(N_i) \geq 2$, the locus
$\Sigma_{N_1,N_2,N_3}$ is at most countable, contained in the
degenerate values $\{0, 1, \infty\}$, the truncation singularity,
and the admissible levels of the parent superalgebra.
\end{proof}


\section{Special cases}%
\label{sec:y-special-cases}

\begin{computation}[$Y_{0,0,N} \simeq \cW_N \times
\mathfrak{gl}(1)$; \ClaimStatusProvedHere]%
\label{comp:y-wn-specialization}
\index{Y-algebra@$Y$-algebra!Y00N@$Y_{0,0,N}$}
With $N_1 = N_2 = 0$, $N_3 = N$, the BRST
definition~\eqref{eq:y-brst-definition} reduces to the
Drinfeld--Sokolov reduction of
$\widehat{\mathfrak{gl}}(N)_\Psi$, giving
$Y_{0,0,N}[\Psi] \simeq \cW_N(k) \times \mathfrak{gl}(1)$
with $k = \Psi - N$.  The strong generators are:
the $\mathfrak{gl}(1)$ current $J$ at spin~$1$ (the
$\mathrm{tr}$ part of $\mathfrak{gl}(N)$), and the
$\cW_N$ generators $W_2 = T, W_3, \ldots, W_N$ at spins
$2, 3, \ldots, N$.

The modular characteristic decomposes as
$\kappa(Y_{0,0,N}) = \kappa(\cW_N) + \kappa(\mathfrak{gl}(1))$.
For $\cW_N$ at level $k = \Psi - N$: the modular characteristic
is computed from the higher-spin self-OPE data.  For
$\mathfrak{gl}(1)$: $\kappa = \Psi$ (the Heisenberg formula).
\end{computation}

\begin{computation}[$Y_{N,0,0} \simeq
\widehat{\mathfrak{gl}}(N)$; \ClaimStatusProvedHere]%
\label{comp:y-affine-specialization}
\index{Y-algebra@$Y$-algebra!YN00@$Y_{N,0,0}$}
By $S_3$ triality (Remark~\ref{rem:y-triality}),
$Y_{N,0,0}[\Psi]$ is the affine $\mathfrak{gl}(N)$ algebra
at level $\Psi' = 1/\Psi$ \textup{(}the M\"obius image under
the triality permutation that sends $(N,0,0)$ to
$(0,0,N)$\textup{)}.  The modular characteristic is the
affine formula:
$\kappa(\widehat{\mathfrak{gl}}(N)_{\Psi'})
= N^2 \cdot (\Psi' + N)/(2N)
= N(\Psi' + N)/2$.
\end{computation}


\section{The genus tower of $Y_{1,1,1}[\Psi]$}%
\label{sec:y111-genus-tower}

\begin{proposition}[{Genus-$1$ free energy of $Y_{1,1,1}[\Psi]$;
\ClaimStatusProvedHere}]%
\label{prop:y111-genus1}
\index{Y-algebra@$Y$-algebra!Y111@$Y_{1,1,1}$!genus tower}
The genus-$1$ free energy is
\begin{equation}\label{eq:y111-genus1}
F_1(Y_{1,1,1}[\Psi])
\;=\; \frac{\Psi}{24}.
\end{equation}
\end{proposition}

\begin{proof}
By genus-$1$ universality
(Theorem~\ref{thm:genus-universality}),
$F_1 = \kappa/24$ unconditionally.  With
$\kappa = \Psi$ (Theorem~\ref{thm:y111-kappa-channels}):
$F_1 = \Psi/24$.
\end{proof}

\begin{proposition}[{Higher-genus tower of $Y_{1,1,1}[\Psi]$;
\ClaimStatusProvedHere}]%
\label{prop:y111-genus-tower}
\index{Y-algebra@$Y$-algebra!Y111@$Y_{1,1,1}$!genus tower}
For all $g \geq 1$, the genus-$g$ free energy of
$Y_{1,1,1}[\Psi]$ is
\begin{equation}\label{eq:y111-genus-tower}
F_g(Y_{1,1,1}[\Psi])
\;=\; \Psi \cdot \lambda_g^{\mathrm{FP}},
\end{equation}
where $\lambda_g^{\mathrm{FP}}$ are the Faber--Pandharipande
intersection numbers on~$\overline{\cM}_g$.  The first values:
\begin{center}
\begin{tabular}{ll}
$F_1 = \Psi/24$, &
$F_2 = 7\Psi/5760$, \\
$F_3 = 31\Psi/967680$, &
$F_4 = 127\Psi/154828800$.
\end{tabular}
\end{center}
\end{proposition}

\begin{proof}
$Y_{1,1,1}[\Psi]$ has two generators: $J$ (weight~$1$)
and $T$ (weight~$2$, $c = 0$).  These have distinct
conformal weights, so the algebra is multi-weight in the
sense of AP27/AP32.  The multi-weight genus expansion
(Theorem~\ref{thm:multi-weight-genus-expansion}) gives
$F_g = \kappa \cdot \lambda_g^{\mathrm{FP}} + \delta F_g^{\mathrm{cross}}$.
The cross-channel correction $\delta F_g^{\mathrm{cross}}$ involves
mixed-propagator graphs between the $J$-channel and the
$T$-channel.  These vanish identically because
$\kappa_T = c/2 = 0$: every mixed graph has at least one
$T$-channel edge, which carries the factor $\kappa_T = 0$.
More precisely, the only mixed OPE $T_{(n)} J$ is the standard
primary OPE (linear, no structure constants beyond the
conformal weight), so the mixed-channel graph amplitudes factor
through $\kappa_T = 0$ at every internal $T$-edge.  Hence
$\delta F_g^{\mathrm{cross}} = 0$ for all $g \geq 2$, and
$F_g = \kappa \cdot \lambda_g^{\mathrm{FP}} = \Psi \cdot
\lambda_g^{\mathrm{FP}}$.
\end{proof}

\begin{remark}[Self-dual coupling]%
\label{rem:y111-self-dual}
\index{Y-algebra@$Y$-algebra!Y111@$Y_{1,1,1}$!self-dual point}
The Koszul duality $\Psi \mapsto -\Psi$ has no real self-dual
fixed point ($\Psi = -\Psi$ gives $\Psi = 0$, which is
outside the domain).  In the complex domain,
$\Psi = 0$ is the degenerate point where $h_1 = 0$ and
the $S_3$ triality degenerates.  The $Y_{1,1,1}$ family has
no self-dual member, unlike the Virasoro ($c = 13$) or
$\cN = 2$ SCA ($c = 3$) families.
\end{remark}


\section{Three-pillar interpretation}%
\label{sec:y-three-pillar}

\begin{remark}[Three-pillar interpretation: $Y_{1,1,1}$]%
\label{rem:y111-three-pillar}
\index{Y-algebra@$Y$-algebra!Y111@$Y_{1,1,1}$!three-pillar}
In the three-pillar architecture
(\S\ref{sec:concordance-three-pillars}):
\begin{enumerate}[label=\textup{(\roman*)}]
\item \emph{Pillar~A (homotopy chiral, MS24):}
  The \v{C}ech complex of $Y_{1,1,1}[\Psi]$ is a non-strict
  homotopy chiral algebra with $F_2 = 0$ (class~$\mathsf{G}$:
  no secondary Borcherds operations at arity~$\geq 3$).  This
  is the Gaussian archetype of Ch$_\infty$-strictification.
\item \emph{Pillar~B (convolution $sL_\infty$, RNW19+Val16):}
  The convolution $sL_\infty$-algebra
  $\operatorname{hom}_\alpha(\cC^!, Y_{1,1,1})$
  is \emph{strictly formal}: all transferred higher brackets
  $\ell_k^{\mathrm{tr}} = 0$ for $k \geq 3$, and the strict dg~Lie
  model is the full $L_\infty$-structure.  This follows from
  class~$\mathsf{G}$ status (Gaussian algebras are $L_\infty$-formal
  at all genera).
\item \emph{Pillar~C (log FM, Mok25):}
  The log FM compactification contributes no planted-forest
  corrections: $\delta_{\mathrm{pf}}^{(g,0)} = 0$ for all
  $g \geq 2$, because the shadow obstruction tower terminates
  at arity~$2$ and all higher-arity codimension-${\geq}2$
  boundary strata contribute zero.
\end{enumerate}
\end{remark}


\section{The holomorphic-topological origin}%
\label{sec:y-ht-origin}

\begin{remark}[Junction of three half-spaces]%
\label{rem:y-junction}
\index{Y-algebra@$Y$-algebra!holomorphic-topological origin}
\index{junction!three half-spaces}
The corner vertex algebra $Y_{N_1,N_2,N_3}[\Psi]$ arises at
the junction of three half-spaces
$\bR_{\geq 0}^{(i)} \times \bC$ ($i = 1,2,3$) meeting along a
common line $\bC$ in the $4$d $\cN = 4$ holomorphic-topological
theory~\cite{gaiotto-rapchak}.  The three gauge groups
$\mathrm{GL}(N_i)$ live on the three half-spaces; the coupling
$\Psi = -\varepsilon_2/\varepsilon_1$ is the ratio of
$\Omega$-background parameters.  The $Y$-algebra is the algebra
of local operators supported on the junction line~$\bC$.

In the holomorphic modular Koszul datum
$\cH(T) = (A, A^!, \mathfrak{C}, r(z), \Theta_A, \nabla^{\mathrm{hol}})$
(\S\ref{sec:concordance-holographic-programme}), the junction
origin provides:
\begin{enumerate}[label=\textup{(\roman*)},nosep]
\item $A = Y_{N_1,N_2,N_3}[\Psi]$,
  $A^! = Y_{N_1,N_2,N_3}[-\Psi]$: the boundary and its
  Koszul dual;
\item $r(z)$: the collision residue
  (Computation~\ref{comp:y111-collision-residue}), encoding the
  OPE singularities along the junction line;
\item $\Theta_A$: the universal MC element, with
  $\kappa = \Psi$ (for $Y_{1,1,1}$) controlling the genus tower;
\item $\nabla^{\mathrm{hol}}$: the holomorphic connection on the
  bundle over $\overline{\cM}_{g,n}$, with curvature
  $\kappa \cdot \omega_g$ at genus~$g \geq 1$.
\end{enumerate}
The $S_3$ triality of the junction descends to the $S_3$
symmetry of the Koszul datum.
\end{remark}
